\section{Construction of the ramified split matching maps}
\label{sec:sr-ramified-matching-construction}

We now construct the finite matching maps at every place of \(T\).
This proves the ramified part of local matching, including its
normalization on arbitrary Schwartz--Bruhat functions.  It does not
supply the archimedean matching quotient, the inert quotient, or the
global comparison of integrals.

Let \(k=F_v\), let \(\mathcal O\) be its valuation ring, and put
\(q=\#(\mathcal O/\varpi)\), \(K_r=\mathrm{GL}_r(\mathcal O)\).
Suppose \(a'<a\), put \(d=a+a'\), and let \(\alpha,\beta\) be
smooth unitary characters of \(k^\times\).  Write
\(\nu=\alpha/\beta\) and assume \(\nu|_{\mathcal O^\times}\ne1\).
The Haar measure on each \(\mathrm{GL}_r(k)\) gives \(K_r\)
mass one.  Set \(X_0=[1_a\ 0]\), \(Y_0=[0\ 1_{a'}]\).
The two split Weil actions, in the designated linear and opposite
linear coordinates, are
\begin{align}
 A(g,h)\phi(X)
 &=\nu(\det g)^{-1}\beta(\det h)
       |\det g|^{-d/2}|\det h|^{-a/2}
                    \phi(g^{-1}Xh^{-t}),\label{eq:sr-ram-action-A}\\
 B(g',h)\phi'(Y)
 &=\nu(\det g')\alpha(\det h)
       |\det g'|^{d/2}|\det h|^{a'/2}
                    \phi'({}^tg'Yh).\label{eq:sr-ram-action-B}
\end{align}
The transpose in the second formula is part of the convention.

\paragraph{Finite support of the ramified Cartan sum.}
For \(r\ge1\), a smooth compactly supported function
\(f:M_r(k)\to\mathbb C\), and \(z\in\mathbb C\), define
\[
 Z_r(f,\nu,z)=
 \sum_{n_1\ge\cdots\ge n_r}
 \int_{K_r\operatorname{diag}(\varpi^{n_i})K_r}
              f(g)\nu(\det g)|\det g|^z\,dg .
\]
This sum has only finitely many nonzero summands.  To prove it,
choose integers \(c,e\) such that \(f=0\) outside
\(\varpi^{-c}M_r(\mathcal O)\), and \(f\) is invariant under
translation by \(\varpi^eM_r(\mathcal O)\).  If \(n_r<-c\),
the corresponding shell does not meet the support: its operator
norm is \(q^{-n_r}\), since multiplication on either side by
\(K_r\) preserves that norm.  Choose
\(u\in\mathcal O^\times\) with \(\nu(u)\ne1\).  On the shell
with diagonal representative \(t_n\), integrate first over
\(k_1,k_2\in K_r\).  The pushforward of their product probability
measure to \(K_rt_nK_r\) is its invariant probability measure.
Replace \(k_1\) by
\(k_1\operatorname{diag}(u,1,\ldots,1)\).  If \(n_1\ge e\),
the argument \(k_1t_nk_2\) changes by an element of
\(\varpi^eM_r(\mathcal O)\), whereas its character factor is
multiplied by \(\nu(u)\).  Thus this shell integral is zero.
Consequently all nonzero shells satisfy
\[
                       -c\le n_r\le n_1<e.
\]
This argument is independent of \(z\).  Each remaining term is
a constant times \(q^{-z\sum_i n_i}\), so \(Z_r\) is an entire
function of \(z\), in fact a Laurent polynomial in \(q^{-z}\).
It also proves linearity directly: take common \(c,e\) for a
finite collection of test functions.

For \(\Re z>r-1\), this sum is the ordinary absolutely convergent
integral over \(\mathrm{GL}_r(k)\).  Indeed additive matrix measure
and multiplicative Haar measure differ by a positive constant
times \(|\det g|^{-r}\).  Successive integral column elimination
on \(M_r(\mathcal O)\) reduces convergence to geometric series
with ratios \(q^{-(\Re z-i)}\), \(0\le i<r\).
A compact support is contained in a scalar dilate of this set.
This proves the assertion without a continuation theorem.
Changes of variables in that half-plane therefore continue to
every \(z\), because both sides are finite Laurent sums.

The same construction applies to \(f(g)=\phi(gX)\) when \(X\)
is any full-row-rank rectangular matrix: choose a right inverse
to \(X\).  It bounds \(g\) on the support, and a sufficiently
small matrix translation of \(g\) changes \(gX\) inside a
translation subgroup for \(\phi\).  In particular every occurrence
of \(Z_a\) below has the finite-support property just proved.

\paragraph{The two maps and their exact common target.}
Define
\begin{align}
 P_A\phi(h)
 &=\beta(\det h)|\det h|^{-a/2}
       Z_a\bigl(g\mapsto\phi(gX_0h^{-t}),\nu,d/2\bigr),
       \label{eq:sr-ram-PA}\\
 P_B\phi'(h)
 &=\alpha(\det h)|\det h|^{a'/2}
       \int_{\mathrm{GL}_{a'}(k)}
          \phi'(gY_0h)\nu(\det g)|\det g|^{d/2}\,dg .
       \label{eq:sr-ram-PB}
\end{align}
The second integral is absolutely convergent since
\(d/2>a'-1\).  Both maps are linear.

Let \(P\) be the upper block parabolic of type \((a,a')\).
The common target is the smooth normalized induced representation
\[
 \Pi_v=\operatorname{Ind}_{P}^{\mathrm{GL}_d(k)}
           \bigl((\alpha\circ\det)\boxtimes(\beta\circ\det)\bigr).
\]
Thus its functions obey, for
\(p=\left(\begin{smallmatrix}A&C\\0&D\end{smallmatrix}\right)\),
\begin{equation}
 F(ph)=\alpha(\det A)\beta(\det D)
       |\det A|^{a'/2}|\det D|^{-a/2}F(h).
 \label{eq:sr-ram-induction-law}
\end{equation}
The exponents are exactly \(\delta_P^{1/2}\).

Here is a verification for both transforms.  One has
\(X_0p^{-t}=A^{-t}X_0\) and \(Y_0p=DY_0\).
In the first integral change \(g\) to \(gA^t\);
its character and determinant factor contribute
\(\nu(\det A)|\det A|^{d/2}\).
Together with the outside factor this is
\[
 \beta(\det A\det D)|\det A\det D|^{-a/2}
        \nu(\det A)|\det A|^{d/2},
\]
which is the factor in \eqref{eq:sr-ram-induction-law}.
For \(P_B\), change \(g\) to \(gD^{-1}\); the resulting factor is
\[
 \alpha(\det A\det D)|\det A\det D|^{a'/2}
        \nu(\det D)^{-1}|\det D|^{-d/2},
\]
the same expression.  The first change of variables is justified
initially for large \(\Re z\), with \(d/2\) replaced by \(z\),
and then by the finite Laurent identity above.

A Schwartz function is fixed by some compact open subgroup
under right matrix multiplication.  Shrink the subgroup to kill
the smooth outside character.  The transforms are consequently
fixed by an open subgroup under right translation, as required
in the smooth induction.  Direct substitution in
\eqref{eq:sr-ram-action-A}--\eqref{eq:sr-ram-action-B} gives
\begin{align}
 P_A(A(g,1)\phi)&=P_A\phi,&
 P_B(B(g',1)\phi')&=P_B\phi',\label{eq:sr-ram-first-invariance}\\
 P_A(A(1,h_1)\phi)(h)&=P_A\phi(hh_1),&
 P_B(B(1,h_1)\phi')(h)&=P_B\phi'(hh_1).
 \label{eq:sr-ram-second-equivariance}
\end{align}
For the first two equalities use respectively the substitutions
\(g_{\rm int}\mapsto gg_{\rm int}\) and
\(g_{\rm int}\mapsto{}^tg'^{-1}g_{\rm int}\).
Their determinant factors cancel the displayed first-group
factors in the actions.  For the latter two equalities use the
representation laws; the first and second actions commute.

Each map is nonzero.  For \(P_A\), choose a principal congruence
subgroup \(K_a(e)\) on which \(\nu\circ\det\) is trivial, and
take \(\phi\) to be the characteristic function of
\[
       [\,K_a(e)\quad \varpi^eM_{a,a'}(\mathcal O)\,].
\]
At \(h=1\) its pullback under \(g\mapsto gX_0\) is precisely
\({\bf1}_{K_a(e)}\), so its transform is
\(\operatorname{vol}(K_a(e))>0\).  This involves only the unit
shell.  The identical construction around \(Y_0\), with the
last square block in \(K_{a'}(e)\), proves \(P_B\ne0\).
If \(a'=0\), replace this last argument by the function \(1\)
on the zero-dimensional matrix space; its transform is the
one-dimensional character \(\alpha\circ\det\).
The applications have \(a'\ge1\).

The one published representation theorem needed at this point
is Bernstein's irreducibility of normalized induction from a
unitary irreducible representation of a Levi subgroup of
\(\mathrm{GL}_d(k)\); see
\href{https://www.math.tau.ac.il/~bernstei/Publication_list/publication_texts/Bernstein-P-invar-SLN.pdf}
{Bernstein, \emph{\(P\)-invariant distributions}, Corollary in
\(\S0.2\), proved in \(\S8.2\)}.
It applies because both inducing representations are the
one-dimensional unitary representations \(\alpha\circ\det\)
and \(\beta\circ\det\).  The induction here is precisely the
normalized one, by \eqref{eq:sr-ram-induction-law}.
No restriction to cuspidal inducing data is being imposed.
Each nonzero invariant image in \(\Pi_v\) is therefore all of
\(\Pi_v\).  Thus \(P_A\) and \(P_B\) are surjective.

\paragraph{Transport to the specified Weil coordinates.}
We record the coordinate argument to specify the model involved.
At a split place,
\(E\otimes_F k=k\oplus k\), with conjugation interchanging the
summands.  Use the designated split lattice bases \(b,c\).
The two complementary tensor summands are Lagrangian; their
symplectic pairing is the matrix trace pairing, with the minus
sign dictated by the skew-Hermitian form.  In linear coordinates
the translation variable is
\[
 b^{-1}u_1c^{-t};
\]
in opposite coordinates it is
\[
                 -{}^tb\,J_Xu_0\,{}^tJ_Wc .
\]
The second formula is also the dual of the first for the
specified symplectic form.  Hence their translations and
modulations give the usual Schr\"odinger representation.
A linear change of coordinates followed, where necessary, by
Fourier transform in one Lagrangian gives isomorphisms
\[
 \iota:\mathcal S_v(V\otimes W)\longrightarrow
                         \mathcal S(M_{a,d}(k)),\qquad
 \iota':\mathcal S_v(V'\otimes W)\longrightarrow
                         \mathcal S(M_{a',d}(k)).
\]
These operations preserve Schwartz--Bruhat functions.

There is a unique lattice-vacuum line in each of these
realizations.  In the matrix model, invariance under integral
translations makes a vector constant on integral cosets;
invariance under integral modulations forces its support to
the annihilator lattice.  When \(\psi_v\) is unramified this is
the integral matrix lattice.  Thus the vacuum is a nonzero
multiple of its indicator.  Scale \(\iota,\iota'\) so that
both multiples equal \(1\).  This is compatible with every
Heisenberg operator.  In a non-self-dual additive convention,
the annihilator lattice must first be scaled by the inverse
different; the places of \(T\) in the present application are
unramified for \(\psi_v\), so that correction is absent.

The restrictions of the splitting characters to the diagonal
\(k^\times\) are trivial at a split place.  Thus their two
components are inverse characters.  For the large pair,
\((\beta/\alpha)(t\circ\sigma)=\nu(t)\) and
\(\beta(t\circ\sigma)=\beta(t)^{-1}\); for the small pair,
\((\alpha/\beta)(t)=\nu(t)\).  The dual-pair actions in these
coordinates preserve the respective Lagrangians.  Egorov's relation makes their operators scalar
multiples of the indicated linear substitutions.  The scalar is
fixed by their action at the origin, namely the Siegel rule with
the specified splitting characters.  For the large pair
\((\beta,\beta/\alpha)\) it gives precisely
\eqref{eq:sr-ram-action-A}; for the small pair
\((\alpha,\alpha/\beta)\) it gives
\eqref{eq:sr-ram-action-B}.  One can check the scalar at any
translated origin as well: translating an arbitrary test vector
reduces the whole equality to its value at zero.  Thus no
additional unit-modulus character can be inserted in these formulas.
The use of the same basis \(c\) for \(W\) ensures that the
second-group variable \(h\) in both formulas is the same
element of \(\mathrm{GL}_d(k)\).  The map from \(U(W)(k)\)
to this general linear group is onto, since its other split
component is uniquely determined by the first and the
skew-Hermitian form.

Define the actual quotient maps by
\[
                         p_v=P_A\iota,\qquad
                         p'_v=P_B\iota'.
\]
They are surjective onto the same nonzero irreducible
\(\Pi_v\).  Equation \eqref{eq:sr-ram-first-invariance} shows
that they annihilate the respective first-group difference
spaces.  Equation \eqref{eq:sr-ram-second-equivariance} gives
exact compatibility under the common second group.
Their matching criterion is pointwise equality
\[
              p_v\phi=p'_v\phi'
        \quad\Longleftrightarrow\quad
              P_A(\iota\phi)=P_B(\iota'\phi')
                    \quad\hbox{on }\mathrm{GL}_d(k).
\]
This is the exact criterion required at a place of \(T\).
Vacuum matching is required only outside \(T\).
The construction is placewise and applies to every allowed
character pair; it makes no choice of a character-dependent
global comparison constant.

\paragraph{Scope for the two rank recursions.}
The preceding construction supplies the ramified finite maps
themselves.  The standard matched vectors and their local cone
evaluations in Section~\ref{sec:ramified-whittaker} consequently
use these very quotient maps.  To infer the cone comparison for
every matched vector one still needs the codimension-two
multiplicity-one statement, together with existence and
equivariance of the descended cone and evaluation operators.
Bernstein irreducibility of \(\Pi_v\) alone does not prove that
branching statement.  Likewise this construction by itself does
not compare the two global theta integrals.  Those are separate
remaining steps in \srcid{ThetaIntegralsAgree}.


\section{Construction at unramified split places}
\label{sec:sr-unramified-split-matching}

We construct the matching maps at every split place outside \(T\).
This includes the possible poles of the unnormalized large-side
matrix integral; evaluation there means the finite normalized
Laurent polynomial constructed below.

Let \(k,\mathcal O,q,\varpi,a,a',d=a+a'\) be as in
Section~\ref{sec:sr-ramified-matching-construction}, with
\(a'<a\).  This time \(\alpha,\beta\), and hence
\(\nu=\alpha/\beta\), are unramified unitary characters.
All multiplicative Haar measures give
\(K_r=\mathrm{GL}_r(\mathcal O)\) mass one.
For a Schwartz--Bruhat function \(f\) on \(M_r(k)\), initially put
\[
 Z_r(f,s)=\int_{\mathrm{GL}_r(k)}
                 f(g)\nu(\det g)|\det g|^s\,dg,
                    \qquad \Re s>r-1 .
\]
Write
\[
 T=\nu(\varpi)q^{-s},\qquad
 D_r(T)=\prod_{i=0}^{r-1}(1-q^iT),\qquad
 c_r=\prod_{j=1}^r(1-q^{-j}).
\]
We prove that
\begin{equation}
                       D_r(T)Z_r(f,s)
                 \quad\hbox{is in }\mathbb C[T,T^{-1}],
 \label{eq:sr-unr-Laurent}
\end{equation}
and that for \(f={\bf1}_{M_r(\mathcal O)}\) this Laurent
polynomial is \(1\).

\paragraph{The integral of the standard lattice.}
Give additive matrix space measure with integral matrices of
mass one.  Then \(dg=c_r^{-1}|\det g|^{-r}\,dX\).
Let
\[
 I_r(t)=\int_{M_r(\mathcal O)}
                         \nu(\det X)|\det X|^t\,dX
\]
on its convergence half-plane; the rank-deficient matrices
have additive measure zero.  In the first column write
\(v=\varpi^j u\), \(u\) primitive, \(j\ge0\).
The corresponding column measure is
\((1-q^{-r})q^{-rj}\).  An element of \(K_r\) sends each
primitive \(u\) to the first basis vector; its determinant is
a unit, so neither character nor absolute-value weights
change.  The remaining determinant is that of an
\((r-1)\)-by-\((r-1)\) integral matrix.  Integration in the
unused first-row coordinates contributes \(1\).
Consequently
\[
 I_r(t)=\frac{1-q^{-r}}
               {1-\nu(\varpi)q^{-(t+r)}}I_{r-1}(t),
                          \qquad I_0(t)=1.
\]
Thus
\[
 I_r(t)=\prod_{j=1}^r
            \frac{1-q^{-j}}{1-\nu(\varpi)q^{-(t+j)}}.
\]
Putting \(t=s-r\) proves
\[
         Z_r({\bf1}_{M_r(\mathcal O)},s)=D_r(T)^{-1}.
\]
It proves absolute convergence for \(\Re s>r-1\) at the same
time.

\paragraph{The polynomial numerator for every test function.}
Here are the details needed to extend the preceding formula
to arbitrary \(f\), rather than merely spherical functions.
After a scalar dilation suppose its support is in
\(M_r(\mathcal O)\), and choose \(e\) such that \(f\) is
constant modulo \(\varpi^eM_r(\mathcal O)\).
We induct on \(r\).
Partition matrices by their residue modulo \(\varpi\).
On a residue class of positive rank \(k\), choose a unit
\(k\)-by-\(k\) pivot block \(A\), after integral row and column
permutations.  The block substitution
\[
 \begin{pmatrix}A&B\\ C&D\end{pmatrix}
       \longmapsto (A,B,C,S),\qquad S=D-CA^{-1}B,
\]
has additive Jacobian \(1\) and determinant
\(\det A\det S\).  The matrix \(A\) is integral invertible,
so its determinant contributes neither character nor norm.
The blocks \(A,B,C\) range over compact sets.
A finite compact-open partition, using the translation level
of \(f\), reduces their integrals to finite linear combinations
of determinant integrals in size \(r-k\), with the same
additive exponent \(t=s-r\).
Their multiplicative parameter is \(s-k\).
By induction their denominator divides
\[
          \prod_{i=k}^{r-1}(1-q^iT),
\]
which divides \(D_r(T)\).
The constants from the additive measures are independent of
\(s\).

The sole zero-residue class is \(X\in\varpi M_r(\mathcal O)\).
Substitute \(X=\varpi X_1\); its scaling factor is \(T^r\).
Repeat the preceding residue partition on this branch.
After at most \(e\) such zero branches, \(f\) is constant
on the remaining central box, and its integral is a scalar
multiple of a dilate of the standard-lattice integral.
This proves \eqref{eq:sr-unr-Laurent}.  The initial scalar
dilation contributes only a Laurent monomial in \(T\).
All operations preserve linearity in \(f\).

It follows that \(Z_r(f,s)\) has the specified rational
continuation and that its normalized value is defined at
every \(s\), even where \(D_r(T)=0\).  No division by a
vanishing value is used in evaluating the normalized map.

\paragraph{The source's finite convolution.}
Expand the convergent Cartan integral as a Laurent series
in an auxiliary variable \(X\), keeping \(s_0=d/2\) fixed:
\[
 \sum_{m\in\mathbb Z}z_mX^m,\qquad
 z_m=\!\!\sum_{\substack{n_1\ge\cdots\ge n_r\\
                         \sum_i n_i=m}}\!
       \int_{K_r\operatorname{diag}(\varpi^{n_i})K_r}
           f(g)\nu(\det g)|\det g|^{s_0}\,dg.
\]
Each inner sum is finite: compact support bounds every
\(n_i\) below, and the sum of the \(n_i\)'s is fixed.
There are only finitely many negative powers in the outer
series.  Multiplying by
\[
 D_{r,s_0}(X)=
       \prod_{i=0}^{r-1}
          (1-\nu(\varpi)q^{\,i-s_0}X)
\]
gives a Laurent polynomial by
\eqref{eq:sr-unr-Laurent}.  Its coefficient of \(X^m\) is
\[
        \sum_{j=0}^r [X^j]D_{r,s_0}(X)\,z_{m-j}.
\]
Only finitely many of these coefficients are nonzero, and
their sum is its value at \(X=1\).  This is exactly the
finite-support convolution called \(\mathrm{zetaUnr}\)
in the source; its polynomial is \(\mathrm{invL}\).
This calculation identifies that definition with the
normalized rational continuation on every Schwartz function.

\paragraph{The common quotient and the spherical vectors.}
Use the actions \(A,B\) and the designated split coordinates
of the preceding section.  At \(s_0=d/2\) put
\begin{align}
 P_A^{\rm unr}\phi(h)
 &=\beta(\det h)|\det h|^{-a/2}
   \left[
     D_a(\nu(\varpi)q^{-s})\,
     Z_a(g\mapsto\phi(gX_0h^{-t}),s)
   \right]_{s=s_0},\label{eq:sr-unr-PA}\\
 P_B^{\rm unr}\phi'(h)
 &=D_{a'}(\nu(\varpi)q^{-s_0})\,
     \alpha(\det h)|\det h|^{a'/2}
       \int_{\mathrm{GL}_{a'}(k)}
       \phi'(gY_0h)\nu(\det g)|\det g|^{s_0}\,dg .
       \label{eq:sr-unr-PB}
\end{align}
The small integral converges ordinarily because
\(s_0>a'-1\).
Moreover its normalizing factor is nonzero:
for \(0\le i<a'\),
\[
             |\nu(\varpi)q^{\,i-s_0}|<1.
\]
The large normalized expression is evaluated as a Laurent
polynomial, as proved above.

Both maps take values in
\[
 \Pi_v=\operatorname{Ind}_{P_{a,a'}}^{\mathrm{GL}_d(k)}
          \bigl((\alpha\circ\det)\boxtimes(\beta\circ\det)\bigr)
\]
with exactly the normalized parabolic law
\eqref{eq:sr-ram-induction-law}.  To verify this, do the same
changes of variables as in the preceding section on a
convergence half-plane, multiply by \(D_r\), and use equality
of rational functions before specializing to \(s_0\).
The identical argument proves first-group invariance and
second-group equivariance.  Smoothness follows from an open
right stabilizer of the input function, as before.

The common target at \(s_0\) must not be confused with a
fixed induction for all parameters.  Before specialization,
the unnormalized large map has parabolic multiplier
\[
 \alpha(\det A)\beta(\det D)
 |\det A|^{a'/2}|\det D|^{-a/2}
                         |\det A|^{\,s-s_0}.
\]
For the first-group action it satisfies
\[
 P_{A,s}(A(g,1)\phi)
                    =|\det g|^{\,s-s_0}P_{A,s}\phi .
\]
These identities follow on the convergence half-plane by the
changes of variables already given; multiply by \(D_a\)
and continue the rational identities.  The normalized
expressions are Laurent polynomials, and the extra scalar
factors are holomorphic and equal to \(1\) at \(s_0\).
Specialization therefore gives the asserted fixed induction
law and first-group invariance, including at a zero of
\(D_a\).  The second-group equivariance has no extra
parameter factor.

The two integral-matrix indicators have value \(1\) at
\(h=1\) under the respective maps, by the standard-lattice
calculation.  They are right \(K_d\)-invariant.  Their images
are equal as functions on all of \(\mathrm{GL}_d(k)\):
Iwasawa decomposition \(G=P_{a,a'}K_d\) and the fixed
parabolic transformation law determine a right-spherical
function from its value at \(1\).  One elementary proof of
this Iwasawa decomposition chooses an integral basis of the
saturated submodule obtained by intersecting the relevant
\(k\)-subspace with \(\mathcal O^d\), then extends it to an
integral basis of \(\mathcal O^d\); the resulting change of
basis lies in \(K_d\) and makes the matrix upper block
triangular.  The common vacuum image is nonzero.

Bernstein's unitary induction theorem applies to this same
\(\Pi_v\), so its two nonzero invariant images are all of
\(\Pi_v\).  Transporting the maps by \(\iota,\iota'\) therefore
gives a common nonzero irreducible quotient of the two
specified local Schwartz spaces, with first-group differences
annihilated, matching preserved by the common second group,
and equal nonzero images of the vacuum vectors.
These are all the finite matching requirements at a split
place outside \(T\).

Together with
Section~\ref{sec:sr-ramified-matching-construction}, this
constructs the finite maps at every split place.  The
construction makes no assertion about the inert places or
the archimedean quotient.  Exact kernel-equals-differences
statements at the additional good places, needed by the
global coinvariant argument, are also separate from the
surjectivity and compatibility proved here.


\paragraph{Zero-rank convention and actual range.}
For \(r=0\), set \(\mathrm{GL}_0(k)=\{1\}\), give it mass
one, set \(D_0=1\), and take the zero-dimensional matrix
Schwartz space to be \(\mathbb C\).
If \(a'=0<a\), the target is the one-dimensional
representation \(\alpha\circ\det\).
The small map is \(c\mapsto(h\mapsto\alpha(\det h)c)\).
The large normalized map has the same covariance under
the whole group \(P_{a,0}=\mathrm{GL}_a(k)\), and sends its
standard lattice vector to this function with \(c=1\).
All assertions follow with these conventions.
In the actual source's admissible triples,
\(2(p+q)+1<a'\), so \(a'\ge2\); the positive-rank proof
above covers every application.


\section{Ramified multiplicity one in codimension two}
\label{sec:sr-ramified-branching}

The irreducibility used to construct the maps in
Section~\ref{sec:sr-ramified-matching-construction} also supplies
the branching statement needed in the local rank recursions.
We prove that statement in matrix coordinates, retaining the
unit diagonal conditions that are essential to its application.

Let \(k\) be a finite extension of \(\mathbb Q_p\),
let \(d=m+2\), and suppose \(n\le m\).
Write a matrix of size \(n\)-by-\(d\) as
\[
                         Y=[\,M\mid x\mid y\,].
\]
Let \(\nu:k^\times\to\mathbb C^\times\) be smooth unitary and
ramified.  Fix \(l\in\mathcal O^\times\) with \(\nu(l)\ne1\).
Let \(P_0\) be the group of matrices
\[
 h=\begin{pmatrix}
       1_m&0&b\\ c&u&e\\0&0&v
    \end{pmatrix},
 \qquad u,v\in\mathcal O^\times .
 \tag{\(*\)}
\]
Here \(c\) is a row of length \(m\), \(b\) a column, and
\(e\in k\).  Its action on the three blocks is
\[
  [M,x,y]h=[\,M+xc,\ ux,\ Mb+ex+vy\,].
\]
For a complex vector space \(W\), call a linear map
\(L:\mathcal S(M_{n,d}(k))\to W\) admissible if
\begin{align}
 L\!\left(\nu(\det a)|\det a|^{d/2}
               \Phi({}^taY)\right)&=L(\Phi)
                       &&(a\in\mathrm{GL}_n(k)),\label{eq:sr-br-left}\\
                         L(\Phi(Yh))&=L(\Phi)
                       &&(h\in P_0).\label{eq:sr-br-right}
\end{align}
These are exactly the two invariances of the source's
\(\mathrm{Inv}\) condition.

Suppose \(W\) is a nonzero countable-dimensional irreducible
module for \(\mathrm{GL}_m(k)\), with action \(\rho\).
Suppose two admissible maps \(L,L'\) satisfy
\begin{equation}
 L(\Phi([Mg,x,y]))=\rho(g)L(\Phi),\qquad
 L'(\Phi([Mg,x,y]))=\rho(g)L'(\Phi)
                      \quad(g\in\mathrm{GL}_m(k)).
 \label{eq:sr-br-middle}
\end{equation}
If \(L\ne0\), then there is \(c\in\mathbb C\) with \(L'=cL\).

If \(n=0\), the matrix Schwartz space is \(\mathbb C\).
A nonzero equivariant \(L\) makes its image a nonzero
\(\rho\)-invariant trivial line in \(W\), so irreducibility
forces \(W\) to be that line.  Then every other equivariant
map is proportional to \(L\).
For the remainder of the proof suppose \(1\le n\le m\).

\paragraph{The local distribution tools.}
We use the localization principle and Frobenius descent for
Schwartz distributions on locally compact totally disconnected
spaces, in
\href{https://www.math.tau.ac.il/~bernstei/Publication_list/publication_texts/Bernstein-P-invar-SLN.pdf}
{Bernstein, \emph{\(P\)-invariant distributions}, \(\S1.4\)--\(\S1.5\)}.
The precise consequence needed is the following.
In a constant-rank matrix chart, let an invariant map to a
Grassmannian record a kernel, and suppose its fibers are
homogeneous spaces for the acting matrix group.
An equivariant distribution on a fiber vanishes if its
stabilizer contains an element on which the equivariance
character differs from the required modular character.
If every fiber has such an element, localization gives
vanishing on the whole chart.  Rank filtrations and the exact
sequence for an open subset and its closed complement then
give the same conclusion for a finite union of these strata.

We check the modular issue explicitly below.  All stabilizer
elements used there are semisimple with eigenvalues powers of
the unit \(l\).  Their adjoint eigenvalues on the corresponding
stabilizer Lie algebra have absolute value \(1\); hence their
modular character is \(1\).  The ambient general linear
groups are unimodular.  The group \(P_0\) is a unipotent group
extended by \((\mathcal O^\times)^2\), whose determinant action
on its unipotent radical has absolute value \(1\), and is
also unimodular.  Thus no modular factor cancels \(\nu(l)\).

For completeness these principles apply to algebraic linear
maps into \(W\).  Compose a map with any linear functional
on \(W\).  On \(\mathcal S(k^N)\), every algebraic functional
is a Schwartz distribution: this space is the inductive limit
of the finite-dimensional spaces of functions supported in a
fixed compact open box and constant modulo a fixed smaller
box.  The same holds for the locally closed matrix strata.
After proving a scalar vanishing statement, separation by the
algebraic dual of \(W\) proves its vector-valued version.

\paragraph{Removal of deficient total rank.}
On matrices of rank \(r<n\), the map \(Y\mapsto\ker Y\) takes
values in \(\operatorname{Gr}_{d-r}(k^d)\); its fibers are the
orbits of left \(\mathrm{GL}_n(k)\).  A nonzero \(n\)-by-\(r\)
minor gives local sections by Gaussian elimination, so this
is precisely a homogeneous-space family of the preceding type.
At a matrix \(Y_0\) of rank \(r<n\), choose a nonzero row
\(c\) with \(cY_0=0\), and a column \(z\) with \(cz=1\).
The projection \(P=zc\) is idempotent and
\[
                a_0=1+(l^{-1}-1)P
\]
fixes \(Y_0\) by left multiplication, has determinant \(l^{-1}\),
and has unit eigenvalues.  Equation \eqref{eq:sr-br-left}, with
\({}^ta=a_0\), supplies the nontrivial multiplier \(\nu(l)\).
Frobenius descent and localization therefore show that the
twisted coinvariants on each deficient-rank stratum are zero.
Induct over the finite rank filtration.  Consequently every
Schwartz function is congruent, modulo the first-group
differences in \eqref{eq:sr-br-left}, to a function compactly
supported on the open full-row-rank locus \(U\).
An admissible \(L\) is determined by its restriction to
\(\mathcal S(U)\).

\paragraph{The first exceptional column.}
On the open set \(x\ne0\), the \(\mathrm{GL}_n(k)\times P_0\)
orbits can be described explicitly.
Reduce modulo the line \(kx\) and put
\[
 r=\operatorname{rank}(\bar M),\qquad
 K=\ker(\bar M:k^m\to k^n/kx).
\]
One further invariant records whether \(y\) lies in
\(\operatorname{im}M+kx\).
For fixed \(r\) and this last binary value, \(K\) ranges over
a Grassmannian and is a complete orbit invariant.
To verify this, first send \(x\) to the first basis vector.
The row \(c\) in \((*)\) removes the first row of \(M\);
left changes of basis on the remaining rows identify matrices
with the same kernel \(K\).  The column \(b\) and entry \(e\)
remove the component of \(y\) in \(\operatorname{im}M+kx\).
A nonzero remaining component can be sent to a fixed last
basis vector by a change of basis on the complement.
On each pivot chart these choices are rational continuous
functions of the nonzero minors.  Thus the invariant map
has homogeneous fibers and the required local sections.

Here is a stabilizer element with nontrivial multiplier at
any such \(Y_0\).  Choose a row \(c_0\) with \(c_0x_0=1\).
Let
\[
 P=x_0c_0,\quad
 Q=e_x(c_0Y_0),\quad
 a_0=1+(l^{-1}-1)P,\quad h_0=1+(l-1)Q,
\]
where \(e_x\) is the column of the exceptional \(x\)-position.
Both \(P,Q\) are idempotent and \(PY_0=Y_0Q\).
It follows that \(a_0Y_0h_0=Y_0\).
The matrix \(h_0\) belongs to \(P_0\): only its \(x\)-row
differs from the identity, its \(x\)-diagonal entry is \(l\),
and its \(y\)-diagonal entry is \(1\).
Its right multiplier is \(1\), while the left multiplier is
\(\nu(l)\), as before.  The modular characters are \(1\).
Localization proves that \(L\) annihilates all functions
supported on \(x\ne0\), and the same is true of \(L'\).

In particular, on \(U\) these maps descend through restriction
to the closed subset \(Z=\{x=0\}\cap U\).  This uses the exact
Schwartz sequence
\[
 0\longrightarrow\mathcal S(U\setminus Z)
 \longrightarrow\mathcal S(U)\longrightarrow\mathcal S(Z)
 \longrightarrow0 .
\]
In this nonarchimedean situation there are no transverse
derivatives of a distribution supported on \(Z\).

\paragraph{The remaining exceptional rank.}
Inside \(Z\), full rank of \(Y=[M,0,y]\) forces either
\(\operatorname{rank}M=n\), or
\(\operatorname{rank}M=n-1\) and \(y\notin\operatorname{im}M\).
The latter is the only exceptional stratum.
Its orbits are again the fibers of \(M\mapsto\ker M\).
At one of its points choose \(c_0M_0=0,\ c_0y_0=1\).
Now put
\[
       P=y_0c_0,\quad
       a_0=1+(l^{-1}-1)P,\quad
       h_0=\operatorname{diag}(1_m,1,l).
\]
Again \(a_0Y_0h_0=Y_0\), \(h_0\in P_0\), and the multiplier
is \(\nu(l)\ne1\), with trivial modular factors.
Gaussian pivot charts and localization prove that the
coinvariants of this exceptional stratum are zero.
It follows that an admissible map is already determined by
its restriction to
\[
                 Z_0=\{\,x=0,\ \operatorname{rank}M=n\,\}.
\]
Precisely, if the induced map on \(\mathcal S(Z_0)\) is zero,
it factors through the exceptional stratum and hence is zero
on \(\mathcal S(Z)\), then on \(\mathcal S(U)\), and finally
on the whole matrix Schwartz space.

\paragraph{Integration in the last column.}
On \(Z_0\), the shear with column \(b\) sends \(y\) to \(y+Mb\).
Since \(M\) is onto, these are all translations in \(y\).
Their coinvariant map is integration in \(y\):
\[
                f(M,y)\longmapsto\int_{k^n}f(M,y)\,dy .
\]
Here additive Haar measure gives \(\mathcal O^n\) mass one.
To check this with varying \(M\), work on a chart where an
\(n\)-by-\(n\) minor \(M_S\) is invertible and use the fiber
coordinate \(z=M_S^{-1}y\).  The shears supported in columns
\(S\) are ordinary translations \(z\mapsto z+t\).
For compactly supported locally constant functions their
coinvariants are Haar integration: subdivide into cosets of
one common compact open subgroup; the zero-sum coefficients
are finite sums of differences of translates.
The Jacobian \(|\det M_S|\) converts integration in \(z\) to
integration in \(y\).  A finite compact-open partition of the
base support joins these chart arguments.

Let
\(\mathcal S_{\rm full}(M_{n,m}(k))\) denote Schwartz functions
compactly supported on the full-row-rank locus.
Choose indicators \(\eta_x,\eta_y\) of \(\mathcal O^n\).
The restriction of \(L\) is therefore encoded by the linear map
\[
 T_L(\psi)
       =L\bigl(\psi(M)\eta_x(x)\eta_y(y)\bigr),
       \qquad \psi\in\mathcal S_{\rm full}(M_{n,m}(k)).
\]
It is independent of the chosen test functions with
\(\eta_x(0)=1\) and \(\int\eta_y=1\).  The previous paragraphs
prove that \(L\mapsto T_L\) is injective.

Equation \eqref{eq:sr-br-left} induces on this lower space
the first-group invariance
\begin{equation}
 T_L\!\left(\nu(\det a)|\det a|^{m/2}
                           \psi({}^taM)\right)=T_L(\psi).
 \label{eq:sr-br-lower-action}
\end{equation}
Indeed \(d/2=(m/2)+1\); changing \(y\) to \({}^ta y\)
in its Haar integral contributes \(|\det a|^{-1}\).
The value at \(x=0\) contributes no determinant.
Equation \eqref{eq:sr-br-middle} gives
\(T_L(\psi(Mg))=\rho(g)T_L(\psi)\), and likewise for \(L'\).

\paragraph{The lower irreducible quotient and scalarity.}
The full-rank matrix space is a principal
\(\mathrm{GL}_n(k)\)-bundle over its Grassmannian of kernels.
Integration along its fibers with weight
\(\nu(\det a)|\det a|^{m/2}\) identifies the quotient in
\eqref{eq:sr-br-lower-action} with
\[
 Q=|\det|^{-n/2}\otimes
       \operatorname{Ind}_{P_{m-n,n}}^{\mathrm{GL}_m(k)}
                    (1\boxtimes(\nu^{-1}\circ\det)).
 \label{eq:sr-br-Q}
\]
The representation on \(Q\) here is pure right pullback.
To check the convention, the integral
\[
 |\det h|^{n/2}\!
 \int_{\mathrm{GL}_n(k)}
       \psi\bigl(a[0\ 1_n]h\bigr)\nu(\det a)|\det a|^{m/2}\,da
\]
has upper-parabolic covariance
\(|\det A|^{n/2}\nu(\det D)^{-1}|\det D|^{-(m-n)/2}\).
Its right action differs from pure pullback by
\(|\det h|^{n/2}\), giving the twist in \eqref{eq:sr-br-Q}.
The integral is compactly supported in its fiber for each
full-rank Schwartz function, so no convergence continuation is
needed here.  On a local trivialization of the frame bundle,
fiber integration has kernel generated by the first-group
differences and is onto: take a compact open fiber bump of
weighted integral \(1\).  A finite partition over the compact
Grassmannian proves the quotient identification.

Both inducing representations in \eqref{eq:sr-br-Q} are
one-dimensional unitary representations.  Bernstein's
irreducibility theorem, with the same normalization as in the
preceding section, makes \(Q\) irreducible; its scalar
\(|\det|^{-n/2}\) twist preserves irreducibility.
When \(m=n\), the parabolic has a zero block and the same
statement is the one-dimensional character directly.

The nonzero map \(T_L:Q\to W\) is an isomorphism: its kernel
and image are invariant, and both modules are irreducible.
Thus \(T_{L'}T_L^{-1}\) is an endomorphism of \(W\) commuting
with \(\rho\).  Such an endomorphism is scalar.
Here is the countable-dimensional form of Schur's argument.
If \(D\) is a nonscalar commuting endomorphism, every
\(D-\lambda\), \(\lambda\in\mathbb C\), is invertible by
irreducibility.  No nonzero polynomial in \(D\) annihilates
a nonzero vector; factoring that polynomial into linear
factors proves this assertion.  Hence for a fixed \(w\ne0\)
the vectors \((D-\lambda)^{-1}w\), \(\lambda\in\mathbb C\),
are linearly independent: a finite relation, multiplied by
the product of its distinct \(D-\lambda\)'s, would give a
nonzero polynomial annihilating \(w\).
This contradicts countable dimension.
Therefore \(T_{L'}=cT_L\), and injectivity of
\(L\mapsto T_L\) proves \(L'=cL\).

\paragraph{Application to the specified local rank steps.}
At a ramified split place the characters on the relevant
opposite model are \(\nu_v\) or \(\nu_v^{-1}\); both are
ramified.  The adapted frame
\(C=[\,\iota\mid H.x\mid H.y\,]\) is invertible and has Gram
matrix \(J_{W_1}\oplus
\left(\begin{smallmatrix}0&1\\-1&0\end{smallmatrix}\right)\).
Use it to order the complementary two columns as \(x,y\).
For \(h\in P_0\) with diagonal entries \(u,v\) on these
columns, its Levi part is represented by a local split
idele \(z\) with components \((u,v^{-1})\).
Its norm is \(1\), and the base-character restriction gives
\[
                \chi_X(z)=\chi_{X,t}(uv)
                          =\chi_{X,t}(\det h).
\]
The source's rank-step equivariance has factor
\(\chi_X(z)|z|^r\), whereas the opposite linear Weil
operator has factor
\(\chi_{X,t}(\det h)|\det h|^{n/2}\).
They agree, so cancellation leaves exactly the pure
right invariance \eqref{eq:sr-br-right}.
The first-group difference condition gives
\eqref{eq:sr-br-left}.

The common lower quotient is \(W=\Pi_{1,v}\).
Let \(\tau\) be its second-group action and identify
\(U(W_1)(k)\) with \(\mathrm{GL}_m(k)\) by the same adapted
frame.  Put
\[
 \rho(g)=
   \bigl(\chi_{X,t}(\det g)|\det g|^{n/2}\bigr)^{-1}
                      \tau(g).
\]
This scalar determinant twist is a representation and
preserves irreducibility.  The same cancellation, with the
middle block \(g\), is precisely
\eqref{eq:sr-br-middle}.  It is irreducible by the finite
matching hypotheses and countable-dimensional because it is
a quotient of the local Schwartz space.  That Schwartz space
has a countable basis obtained from compact open boxes and
their cosets.  The rank inequalities in the two steps give
\(n\le d_1=m\), as required above.

Thus two nonzero descended cone and evaluation maps of this
specified type are proportional.  Evaluation on the standard
matched vectors computed in
Section~\ref{sec:ramified-whittaker} fixes the proportionality
constant to \(\epsilon_v(-\delta)\) in the beta step, and to
\(\epsilon_v(1-\delta)\) in the alpha step.
This proves the ramified multiplicity-one input used for that
extension from the standard vectors.  Existence of the
descended operators and the global constant-term argument
are not consequences of multiplicity one and remain separate
requirements in the full rank recursion.


\section{Good split places and the exact coinvariant condition}
\label{sec:sr-good-split-coinvariants}

The finite matching constructions imply the stronger kernel
statement needed at the source's good places.  We give both
ingredients: nontriviality of the local character, and the
identification of the full Schwartz coinvariant quotient.
This proves the \(\mathrm{Coinv}\) component of
\srcid{ThetaIntegralsAgree}, with its original quantifiers
over any compatible matching maps.

\paragraph{Why the local character is nontrivial.}
Let \(E/\mathbb Q\) be Galois and CM, with CM type \(\Sigma\).
Let \(\nu\) be a unitary Hecke character whose infinity type is
\[
        \nu_\infty(z)=
          \prod_{\tau\in\Sigma}
             \left(\frac{\tau z}{|\tau z|}\right)^{\kappa_\tau},
                       \qquad \kappa_\tau>0.
\]
Suppose \(\nu\) is unramified away from a finite set \(T\).
Let a rational prime \(p\) split completely in \(E\), and
let \(P\mid p\) be outside \(T\).  Then the local
character at \(P\) is nontrivial.

Indeed assume that the character at \(P\) is trivial.
Choose \(h>0\) with \(P^h=(\pi)\), \(\pi\in\mathcal O_E\);
the ideal class group is finite.  At every other finite
prime \(\pi\) is a unit.  Its character value is \(1\) outside
\(T\); at the finitely many primes of \(T\), its values are
roots of unity.  For the latter assertion, the continuous
character of the compact local unit group has finite image:
it is a compact totally disconnected subgroup of
\(\mathbb C^\times\), hence finite.  Choose \(M>0\) killing
all these values.  The finite-idele value at \(\pi^M\) is
therefore \(1\), and the Hecke relation implies
\[
        \prod_{\tau\in\Sigma}
          \left(\frac{\tau\pi}{|\tau\pi|}\right)^{M\kappa_\tau}
                          =1.
\]
Squaring and using
\(|\tau\pi|^2=\tau\pi\,\overline{\tau\pi}\) gives
\begin{equation}
 \prod_{\tau\in\Sigma}(\tau\pi)^{M\kappa_\tau}
     =\prod_{\tau\in\Sigma}(\overline{\tau\pi})^{M\kappa_\tau}.
 \label{eq:sr-good-character-contradiction}
\end{equation}

Fix \(\tau_0\in\Sigma\).  Since \(E/\mathbb Q\) is Galois,
write \(\tau=\tau_0g_\tau\), \(g_\tau\in\operatorname{Gal}(E/\mathbb Q)\),
with \(g_{\tau_0}=1\), and let \(c\) be complex conjugation.
Pull \eqref{eq:sr-good-character-contradiction} back to \(E\)
and take its \(P\)-adic valuation.  The left side has a
strictly positive contribution from \(g_{\tau_0}=1\).
The right side has none: a factor could have positive
valuation only if \(g_\tau c\) fixed \(P\).  A completely split
prime has trivial decomposition group, so this would force
\(g_\tau=c\), that is \(\tau=\bar\tau_0\), which is not in the
CM type \(\Sigma\).  This contradiction proves the assertion.
The same argument applies to every \(P\mid p\) outside \(T\).

In the present packet \(\nu=\alpha/\beta\) and
\(\kappa_\tau=m_\alpha(\tau)-m_\beta(\tau)>d>0\).
The source's \(\mathrm{Good}\) condition says that \(v\notin T\)
and its underlying rational prime splits completely in \(E\).
Thus every designated local component \(\nu_t\), for either
embedding \(t:E\to F_v\), is nontrivial.  Such a place splits
in \(E/F\), since complete splitting gives local degree one
in every intermediate field.

\paragraph{The full-rank quotient.}
Let \(k\) be nonarchimedean of characteristic zero, let
\(1\le n\le d\), and let \(\nu\ne1\) be smooth unitary.
Consider either of the two matrix actions
\[
 \nu(\det g)^{-1}|\det g|^{-d/2}\phi(g^{-1}X),
 \qquad
 \nu(\det g)|\det g|^{d/2}\phi({}^tgX).
\]
The additional character and determinant factor for the
second group does not change the first-group differences.
Let \(\mathcal R\) be the span of those differences, and
\(\mathcal S_{\rm full}\) the functions compactly supported
on full-row-rank matrices.

Every class in \(\mathcal S(M_{n,d}(k))/\mathcal R\)
has a representative in \(\mathcal S_{\rm full}\).
To see this, filter the closed deficient-rank locus by rank.
On its rank-\(r\) stratum, the kernel map to
\(\operatorname{Gr}_{d-r}(k^d)\) has homogeneous fibers under
\(\mathrm{GL}_n(k)\); Gaussian pivot charts give local
sections.  The stabilizer of
\(\left(\begin{smallmatrix}1_r&0\\0&0\end{smallmatrix}\right)\)
contains
\[
                \begin{pmatrix}1_r&B\\0&D\end{pmatrix},
              \qquad D\in\mathrm{GL}_{n-r}(k).
\]
Its determinant ranges over all of \(k^\times\).
The modular character is a positive real power of
\(|\det D|\), as is the modulus in either displayed action.
For equivariant distributions on this homogeneous space to
exist, their character restricted to this stabilizer must
equal that positive modular character.
This is impossible: choose \(D\) with
\(\nu(\det D)\ne1\).  A nontrivial unit-modulus number cannot
be a positive real number.  Frobenius descent and Bernstein
localization, with the kernel map just specified, give zero
coinvariants on each deficient-rank stratum.
The exact Schwartz sequences for the rank filtration show
that the quotient of the whole Schwartz space by
\(\mathcal S_{\rm full}\) has zero coinvariants.
Right exactness of coinvariants proves the claimed
surjectivity from \(\mathcal S_{\rm full}\).

The full-row-rank frame bundle over the Grassmannian has
free left \(\mathrm{GL}_n(k)\)-fibers.
Fiber integration with the displayed character and norm
weight identifies its coinvariants with a normalized
parabolic induction of two unitary determinant characters,
up to a scalar determinant twist for the right action.
The verification is the same fiber-bump and finite
partition argument as for
\eqref{eq:sr-br-Q}.  For the two complementary matrix models
with \(n=a\) or \(a'\), it is the common induction
\[
 \operatorname{Ind}_{P_{a,a'}}^{\mathrm{GL}_d(k)}
          ((\alpha\circ\det)\boxtimes(\beta\circ\det))
\]
in the action convention of their matching maps.
It is irreducible by Bernstein's unitary induction theorem.

Consequently the natural surjection from the full-rank
coinvariant quotient onto the whole Schwartz coinvariant
quotient is either zero or an isomorphism.  It is nonzero.
At a ramified split place the explicit full-rank bump in
Section~\ref{sec:sr-ramified-matching-construction} has
nonzero transform.  At an unramified split place with
\(\nu\ne1\), the same bump works for the normalized maps:
the large-side normalizing factor does not vanish, because
a zero of
\(\prod_{i<a}(1-\nu(\varpi)q^{\,i-d/2})\)
would require \(i=d/2\) and \(\nu(\varpi)=1\).
An unramified nontrivial character has
\(\nu(\varpi)\ne1\).  The small-side factor is always
nonzero by \(a'<d/2\).
Thus the whole coinvariant quotient is itself this nonzero
irreducible representation.

This argument uses right exactness only.  It does not
assume that coinvariants preserve injections.

\paragraph{Any matching map has the exact kernel.}
Now let \(p\) be any of the source's finite matching maps:
it is onto a nonzero irreducible \(\Pi\) and annihilates
\(\mathcal R\).
The induced map
\[
       \mathcal S(M_{n,d}(k))/\mathcal R\longrightarrow\Pi
\]
is a nonzero surjection between irreducible modules, hence
an isomorphism.  Therefore, on Schwartz functions,
\[
                       p\phi=0\quad\Longleftrightarrow\quad
                                      \phi\in\mathcal R.
\]
Moreover \(\mathcal S_{\rm full}\) surjects under \(p\),
because it already surjects onto the whole coinvariant
quotient.  Both assertions survive the linear and opposite
Weil coordinate isomorphisms.

If \(\ell\) is any linear functional on the ambient
function space that annihilates \(\mathcal R\), its
restriction to Schwartz functions factors uniquely through
\(p\).  Define \(\ell_0(p\phi)=\ell(\phi)\); the kernel
identity proves well-definedness and surjectivity of \(p\)
gives its entire domain.  This is exactly the functional
factorization demanded in \(\mathrm{Coinv}\), rather than
a statement about merely continuous functionals.

Apply this argument at a good \(v\) using either embedding
\(t:E\to F_v\).  The corresponding opposite linear
realization sends the source's
\(\mathcal S^\circ_{\mathrm{loc}}(t)\) to the full-row-rank
Schwartz space.  Since the local character is nontrivial
for each \(t\), its image under \(p'_v\) is all of \(\Pi_v\).
Changing the split embedding may change that full-rank
subspace, but the surjectivity argument applies separately
to each one and uses the same already chosen quotient map.

\paragraph{Stability of the global small subspace.}
At the distinguished split place \(v_1\), the opposite linear
Weil action consists of invertible left and right matrix
substitutions times a scalar, as in
\eqref{eq:sr-ram-action-B}.  Such substitutions preserve
matrix rank.  The local full-rank subspace is therefore
stable under both members of the dual pair.
At all other finite places the local actions preserve
Schwartz--Bruhat functions; at almost all places the
adelic group element is integral and fixes the vacuum.
At infinity the Weil action preserves Schwartz functions.
The restricted tensor product action consequently carries
each defining finite sum of pure tensors for
\(\mathcal S^\circ\) to another such finite sum, with its
\(v_1\)-factor still in the same full-rank subspace.
Thus \(\omega'(\gamma)\mathcal S^\circ
\subseteq\mathcal S^\circ\) for every adelic pair \(\gamma\).

Combining the exact kernels, the full-rank surjections for
every split embedding, the factorization of every
annihilating functional, and this global stability proves
all clauses of \(\mathrm{Coinv}\).
It does not supply the nonzero global theta integral or the
two character-uniform rank-step constants in
\srcid{ThetaIntegralsAgree}.
